\section*{Capítulo \ref{ch:b2:atmospheric-thermodynamics} --- Termodinámica atmosférica}

\begin{solution}{exo:b2:atmospheric-thermodynamics:1}
(a) $e = e_s(\qty{278.15}{K}) = \qty{8.73}{hPa}$ y $e_s(\qty{283.15}{K}) = \qty{12.32}{hPa}$.
(b) $8.73/12.32 = \qty{70.9}{\percent}$. (c) $r = 0.622 \times 8.73/(850 - 8.73) =
\qty{6.45}{g/kg}$ y $q = r/(1 + r) = \qty{6.41}{g/kg}$.
\end{solution}

\begin{solution}{exo:b2:atmospheric-thermodynamics:2}
(a) $\mathrm{d}\ln e_s/\mathrm{d}T = 5419/T^2$: un \qty{7.26}{\percent} por kelvin a
\qty{273.15}{K}, un \qty{6.31}{\percent} a \qty{293.15}{K} y un \qty{5.90}{\percent} a
\qty{303.15}{K}. (b) $e_s(\qty{20}{\celsius}) = \qty{23.66}{hPa}$ y $e_s(\qty{25}{\celsius})
= \qty{32.26}{hPa}$; cociente 1.363, frente a $1.063^5 = 1.357$. La tasa es una derivada de
$\ln e_s$, así que en \qty{5}{K} el factor es $\exp(5 \times 0.0620) = 1.363$, con 0.0620 la
tasa en el punto medio, \qty{22.5}{\celsius}; componer un \qty{6.3}{\percent} cinco veces se acerca
pero no es exacto. (c) $23.66/23.39 = 1.012$: la fórmula da un \qty{1.2}{\percent} de más, porque
conserva el calor latente del punto triple, mientras que el calor latente disminuye cuando la temperatura
sube; la pendiente de la fórmula es demasiado fuerte, y su error crece con la distancia a
\qty{0}{\celsius}.
\end{solution}

\begin{solution}{exo:b2:atmospheric-thermodynamics:3}
(a) $125 \times (30 - 14) = \qty{2000}{m}$. (b) $30 - 9.8 \times 2.0 = \qty{10.4}{\celsius}$.
(c) $14 - 1.8 \times 2.0 = \qty{10.4}{\celsius}$: la temperatura y el punto de rocío se encuentran, como deben hacerlo en
el nivel de saturación; la regla es este cálculo. (d) $125 \times (33 - 14) = \qty{2375}{m}$:
las bases suben a medida que la tarde calienta el aire próximo al suelo.
\end{solution}

\begin{solution}{exo:b2:atmospheric-thermodynamics:4}
(a) La curva del punto de rocío coincide con la de la temperatura desde el suelo hasta unos \qty{908}{hPa}:
el medio kilómetro inferior está saturado, o casi. (b) Entre unos 908 y \qty{889}{hPa}
la curva de la temperatura gira hacia la derecha, hacia valores más cálidos, cuando la presión disminuye: la
temperatura aumenta con la altura. (c) El NCA, hacia \qty{964}{hPa}, unas decenas de metros más arriba; el NCL, hacia
\qty{730}{hPa}; el NE, hacia \qty{235}{hPa}. (d) En la cima de la inversión el punto de rocío
está \qty{11}{K} por debajo de la temperatura, y \qty{20}{K} por debajo a \qty{850}{hPa}: un aire cálido y muy seco
descansa sobre un aire fresco y húmedo. Su humedad relativa es baja, así que no adquirió su
calor por condensación; el aire que desciende se calienta por compresión según el gradiente adiabático seco
conservando su \omterm{def:b2:atmospheric-thermodynamics:moist-air}{razón de mezcla}, lo que daría una capa así.
\end{solution}

\begin{solution}{exo:b2:atmospheric-thermodynamics:5}
(a) $T_v = 303.15 \times (1 + 0.020/0.622)/1.020 = \qty{306.76}{K}$; la aproximación da
$303.15 \times (1 + 0.61 \times 0.020) = \qty{306.85}{K}$. (b) $\rho = p/(R_dT_v) = 10^5/(287.0
\times 306.76) = \qty{1.1358}{kg/m^3}$, frente a $10^5/(287.0 \times 303.15) =
\qty{1.1494}{kg/m^3}$ para el aire seco: un \qty{1.2}{\percent} más ligero. (c) En $T_v - T =
\qty{3.6}{K}$. (d) Una porción que asciende solo necesita ser más ligera que su entorno: sobre mares cálidos
la humedad del aire añade una flotación equivalente a varios kelvin, tanto como un exceso
de temperatura considerable.
\end{solution}

\begin{solution}{exo:b2:atmospheric-thermodynamics:6}
(a) $e_s$ = 6.11, 23.66 y \qty{43.54}{hPa}; $r_s = 0.622\,e_s/(1000 - e_s)$ = 3.82, 15.07 y
\qty{28.31}{g/kg}; $\Gamma_s$ = 6.48, 4.20 y \qty{3.44}{K/km}. (b) La porción seca se enfría
\qty{9.8}{K}, la saturada \qty{4.2}{K}: la porción saturada está \qty{5.6}{K} más cálida.
(c) Más arriba el aire es más frío, $r_s$ es menor, queda poco vapor por condensar y
$\Gamma_s$ vuelve hacia $\Gamma_d$.
\end{solution}

\begin{solution}{exo:b2:atmospheric-thermodynamics:7}
(a) $e = e_s(\qty{291.65}{K}) = \qty{21.52}{hPa}$, $r = 0.622 \times 21.52/(971.1 - 21.52) =
\qty{14.1}{g/kg}$; $\theta = 292.15 \times (1000/971.1)^{0.2857} = \qty{294.6}{K}$; $\theta_e =
294.6 \exp[2.501 \times 10^6 \times 0.0141/(1004.5 \times 291.5)] = 294.6 \times 1.128 =
\qty{332.3}{K}$. (b) $e_s(\qty{264.25}{K}) = \qty{3.13}{hPa}$, $r_s = 0.622 \times 3.13/(500 -
3.13) = \qty{3.92}{g/kg}$; $\theta = 264.25 \times 2^{0.2857} = \qty{322.1}{K}$; la porción está
saturada, así que $T_{\text{NCA}} = T$ y $\theta_e = 322.1 \exp[2.501 \times 10^6 \times
0.00392/(1004.5 \times 264.25)] = \qty{334.3}{K}$. (c) $\theta$ aumentó \qty{27.5}{K}, calentada
por el calor latente de los \qty{10.2}{g/kg} de vapor que se condensaron; $\theta_e$ cambió un
\qty{0.6}{\percent}, dentro de la precisión de su fórmula aproximada: se conserva.
\end{solution}

\begin{solution}{exo:b2:atmospheric-thermodynamics:8}
(a) $\text{CAPE} = 287.0 \times 4 \times \ln(700/300) = 287.0 \times 4 \times 0.847 =
\qty{973}{J/kg}$. (b) $\sqrt{2 \times 973} = \qty{44.1}{m/s}$. (c) Entre 1 y
\qty{1500}{J/kg}: positiva, pero no grande. (d) Una banda de anchura horizontal constante entre dos
curvas; por la \cref{prop:b2:atmospheric-thermodynamics:areas}, su área es proporcional a $\int
\Delta T\,\mathrm{d}\ln p = 4 \ln(700/300)$, sea cual sea la inclinación de las isotermas.
\end{solution}

\begin{solution}{exo:b2:atmospheric-thermodynamics:9}
(a) $D_c = \sqrt{3 \times 0.61 \times 80^3/2.1}\,\unit{nm} = \qty{668}{nm} =
\qty{0.668}{\micro m}$; $s_c = 2A/(3D_c) = 2 \times 2.1/(3 \times 668) = 0.00210 =
\qty{0.210}{\percent}$. (b) Al resolver $s_c = \sqrt{4A^3/(27\kappa D_d^3)}$ respecto de $D_d$: $D_d =
\bigl(4A^3/(27\kappa s_c^2)\bigr)^{1/3}$, \qty{63}{nm} al \qty{0.3}{\percent} y \qty{131}{nm} al
\qty{0.1}{\percent}. (c) Todas las partículas de entre 63 y \qty{131}{nm} se activan también: más
gotitas y, para la misma agua líquida, más pequeñas.
\end{solution}

\begin{solution}{exo:b2:atmospheric-thermodynamics:10}
(a) $w = \sqrt{2 \times 130} = \qty{16.1}{m/s}$, imposible para un aire en reposo cerca del suelo.
(b) La temperatura potencial de la cima de la inversión es $293.65 \times (1000/889.1)^{0.2857} =
\qty{303.7}{K}$; una porción de superficie con esta $\theta$ tiene $T = 303.7 \times (971.1/1000)^{0.2857}
= \qty{301.2}{K}$, \qty{28.0}{\celsius}. (c) $125 \times (28.0 - 18.5) = \qty{1190}{m}$ sobre el
suelo, por encima de la antigua cima de la inversión. (d) Un frente, que eleva el aire sobre su pendiente; el lento ascenso
del aire que precede a una perturbación en altura, que enfría la propia tapa.
\end{solution}

\begin{solution}{exo:b2:atmospheric-thermodynamics:11}
(a) Cuando $r_s \to 0$, el numerador tiende a $g$ y el denominador a $c_p$. (b) Escríbase
$L^2r_s/(R_vT^2) = c_p \times \bigl(Lr_s/(R_dT)\bigr) \times LR_d/(R_vTc_p)$, con
$LR_d/(R_vTc_p) = 5.1$ a \qty{303}{K}: cuando $Lr_s/(R_dT) \gg 1$, también $L^2r_s/(R_vT^2) \gg c_p$,
y $\Gamma_s \to g\,(Lr_s/(R_dT))/(L^2r_s/(R_vT^2)) = gR_vT/(LR_d) = 9.81 \times 461.5 \times
303/(2.501 \times 10^6 \times 287.0) = \qty{1.91}{K/km}$. (c) Con $T_d = T$,
$\Gamma_{\text{roc}} = gR_vT/(LR_d)$, el mismo. En ese límite el calor latente es tan grande que
la porción apenas puede condensar nada de su vapor sin calentarse muy por encima de su punto de
saturación: conserva su \omterm{def:b2:atmospheric-thermodynamics:moist-air}{razón de mezcla} y sigue la línea de \omterm{def:b2:atmospheric-thermodynamics:moist-air}{razón de mezcla} constante, a lo largo de la cual la
temperatura baja según $\Gamma_{\text{roc}}$. (d) A \qty{30}{\celsius} y \qty{1000}{hPa},
$Lr_s/(R_dT) = 0.81$, no mucho mayor que 1, y $L^2r_s/(R_vT^2) = 4180$, cuatro veces $c_p$:
$\Gamma_s = \qty{3.4}{K/km}$, aún lejos de \qty{1.9}{K/km}. El límite no se alcanza.
\end{solution}

\begin{solution}{exo:b2:atmospheric-thermodynamics:12}
(a) $2.868/2.599 = 1.104$ y $1.257/1.032 = 1.218$: un \qty{10.4}{\percent} y un
\qty{21.8}{\percent}. (b) Las gotitas se activan con unas décimas de uno por ciento de \omterm{def:b2:atmospheric-thermodynamics:microphysics}{sobresaturación},
decenas de veces menos. (c) Los cristales crecen deprisa por deposición; toman vapor del aire,
que queda por debajo de la saturación respecto de las gotitas; las gotitas se evaporan y su vapor pasa al
hielo (el proceso de Bergeron), hasta que la nube se glacia. (d) A \qty{0}{\celsius} las dos
presiones de vapor de saturación son iguales: se encuentran en el punto triple, y el cociente tiende a uno
cerca de él.
\end{solution}

\begin{solution}{pb:b2:atmospheric-thermodynamics:1}
\textbf{1.} $e = e_s(\qty{291.65}{K}) = \qty{21.5}{hPa}$.
\textbf{2.} $r = 0.622 \times 21.5/(971.1 - 21.5) = \qty{14.1}{g/kg}$.
\textbf{3.} $q = 14.1/1.0141 = \qty{13.9}{g/kg}$.
\textbf{4.} $e_s(\qty{292.15}{K}) = \qty{22.2}{hPa}$; $21.5/22.2 = \qty{97}{\percent}$.
\textbf{5.} $0.61 \times 0.0141 \times 292.15 = \qty{2.5}{K}$.
\textbf{6.} $125 \times 0.5 = \qty{63}{m}$; $p = 971.1 \exp[-9.81 \times 62.5/(287.0 \times 294.7)]
= \qty{964}{hPa}$.
\textbf{7.} $\theta = 292.15 \times (1000/971.1)^{0.2857} = \qty{294.6}{K}$.
\textbf{8.} $16.9 - 17.3 = \qty{-0.4}{K}$: más fría.
\textbf{9.} $(20.5 - 16.4)/(1102 - 924) = 4.1/178$: la temperatura sube \qty{23}{K/km}, un
gradiente de \qty{-23}{K/km}. Una inversión: absolutamente estable, una tapa.
\textbf{10.} $15.4 - 20.5 = \qty{-5.1}{K}$.
\textbf{11.} $e_s(\qty{-1.8}{\celsius}) = \qty{5.4}{hPa}$ y $e_s(\qty{18.6}{\celsius}) =
\qty{21.7}{hPa}$: un \qty{25}{\percent}. Cálido y seco, puede ser aire que descendió y se calentó por
compresión, o aire calentado sobre suelo seco a barlovento; no es el aire húmedo de debajo.
\textbf{12.} $18.6 + 18.4 = \qty{37.0}{K}$ en $5725 - 1489 = \qty{4236}{m}$, \qty{8.73}{K/km}: entre $\Gamma_s$ y
$\Gamma_d$, condicionalmente inestable.
\textbf{13.} $-18.4 - (-8.9) = \qty{-9.5}{K}$: inestabilidad extrema (de $-8$ a $-10$).
\textbf{14.} Desde el suelo hacia arriba: 0, $-0.4$, $-0.2$, $-5.1$, $-4.9$, $+1.6$, $+9.5$, $+10.0$,
$+5.6$, $+1.5$, $-9.5$~\unit{K}.
\textbf{15.} Entre 850 ($-4.9$) y \qty{700}{hPa} ($+1.6$), en la fracción $4.9/6.5 = 0.754$:
$\ln p = \ln 850 + 0.754 \ln(700/850)$, $p = \qty{734}{hPa}$; $z = 1489 + 0.754 \times 1621 =
\qty{2711}{m}$, \qty{2.37}{km} sobre el suelo.
\textbf{16.} Entre 250 ($+1.5$) y \qty{200}{hPa} ($-9.5$), en la fracción $1.5/11.0 =
0.136$: $p = \qty{242.5}{hPa}$ y $z = 10560 + 0.136 \times 1446 = \qty{10757}{m}$, unos
\qty{200}{m} por debajo de la tropopausa señalada (\qty{10954}{m}).
\textbf{17.} Los tramos ($\Delta T$ medio por el $\ln$ del cociente de presiones): $0.80 \times 0.0474
= 0.038$; $5.55 \times 0.3365 = 1.867$; $9.75 \times 0.2231 = 2.176$; $7.80 \times 0.2877 =
2.244$; $3.55 \times 0.1823 = 0.647$; $0.75 \times 0.0305 = 0.023$; suma \qty{6.995}{K}.
$\text{CAPE} = 287.0 \times 6.995 = \qty{2010}{J/kg}$, con tres cifras.
\textbf{18.} Entre 1500 y \qty{2500}{J/kg}: grande.
\textbf{19.} $-0.2 \times 0.0486 - 0.3 \times 0.0189 - 2.65 \times 0.0207 - 5.0 \times 0.0450 -
2.45 \times 0.1467 = -0.655$~\unit{K}; $\text{CIN} = 287.0 \times (-0.655) = \qty{-188}{J/kg}$.
\textbf{20.} La tabla solo conserva once niveles: entre ellos, la regla de los trapecios sustituye las
curvas por segmentos rectos en $\ln p$ y pierde sus inflexiones. La malla de \qty{1}{hPa} sigue las curvas. La \omterm{def:b2:atmospheric-thermodynamics:virtual}{temperatura virtual}
añade la flotación del vapor de agua de la porción, más abundante en la porción saturada que en
el aire seco que la rodea: una CAPE mayor, una CIN menor.
\textbf{21.} $\sqrt{2 \times 2010} = \qty{63}{m/s}$.
\textbf{22.} $10757 - 2711 = \qty{8046}{m}$ a $63.4/2 = \qty{31.7}{m/s}$: \qty{254}{s}, unos cuatro
minutos.
\textbf{23.} El peso del agua condensada que transporta la corriente ascendente; la mezcla con el aire más seco
que la rodea (el arrastre), que la enfría al evaporarse las gotitas.
\textbf{24.} La teoría de la porción da a la corriente ascendente de una tormenta alimentada por este aire una velocidad de unos
\textbf{\qty[text-series-to-math=true]{63}{m.s^{-1}}}, una cota superior.
\end{solution}
